% macshell.tex - Mac-style terminal environment definition

\usepackage{xcolor}
\usepackage{tcolorbox}
\usepackage{listings}
\usepackage{etoolbox}

% Load tcolorbox libraries
\tcbuselibrary{listings,skins}

% Define macOS window button colors
\definecolor{macred}{HTML}{FF5F57}
\definecolor{macyellow}{HTML}{FEBC2E}
\definecolor{macgreen}{HTML}{28C840}

% Atom One Dark Colors
\definecolor{atom-bg}{HTML}{282c34}   
\definecolor{atom-fg}{HTML}{abb2bf}   
\definecolor{atom-comment}{HTML}{5c6370} 
\definecolor{atom-keyword}{HTML}{c678dd} 
\definecolor{atom-string}{HTML}{9c6f35}  
\definecolor{atom-number}{HTML}{caf291}  
\definecolor{atom-type}{HTML}{3beb4f}
\definecolor{atom-notation}{HTML}{b8b848}

% Define Kotlin Language for listings
\lstdefinelanguage{Kotlin}{
  comment=[l]{//},
  morecomment=[s]{/*}{*/},
  morestring=[b]",
  morestring=[s]{"""*}{*"""},
  morestring=[s]{'''*}{*'''},
  alsoletter={@}, % Treat @ as part of keywords for annotations
  morekeywords=[1]{
    % Kotlin standard keywords
    abstract, actual, annotation, as, break, by, catch, class, companion,
    const, constructor, continue, crossinline, data, delegate, do, dynamic,
    else, enum, expect, external, false, final, finally, for, fun, get, if,
    import, in, infix, init, inline, inner, interface, internal, is, it,
    lateinit, noinline, null, object, open, operator, out, override, package,
    private, protected, public, reified, return, sealed, set, super, suspend,
    tailrec, this, throw, true, try, typealias, typeof, val, var, vararg,
    when, where, while,
  },
  morekeywords=[2]{
    % Data Types / Classes (Colored by your custom green: 3beb4f)
    String, Int, Boolean, Float, Double, Long, Short, Byte, Char,
    Array, List, Set, Map, Unit, Any, Nothing, Collection, Iterable,
    % You can add your own custom classes here too!
    UIState, User, ViewModel, Context, Modifier
  },
  morekeywords=[3]{
    @Composable, @Preview, @OptIn, @JvmStatic, @JvmOverloads
  },
  sensitive=true
}

% Define the Listing Style
\lstdefinestyle{atom-one-dark}{
  backgroundcolor  = \color{atom-bg},
  basicstyle       = \ttfamily\color{atom-fg}\scriptsize,
  commentstyle     = \color{atom-comment},
  keywordstyle     = [1]\color{atom-keyword},  % [1] Standard Keywords (Purple)
  keywordstyle     = [2]\color{atom-type},     % [2] Data Types & Classes (Green)
  keywordstyle     = [3]\color{atom-notation},
  stringstyle      = \color{atom-string},
  numberstyle      = \color{atom-number},
  frame            = none, 
  showstringspaces = false,
  breaklines       = true,
  columns          = fullflexible,
  keepspaces       = true,
}

% USE \newtcblisting instead of \newtcolorbox for code blocks!
\newtcblisting{macshell}[1][]{
  enhanced,
  colback=atom-bg,
  colframe=atom-bg,
  boxrule=0pt,
  arc=3mm,
  top=20pt,     
  bottom=5pt,
  left=10pt,
  right=10pt,
  before skip=10pt,
  after skip=10pt,
  code={\setLTR}, % <-- MAGIC 1: Forces the box itself to be explicitly Left-to-Right
  overlay={
    % Draw window control buttons
    \node[circle, fill=macred,   minimum size=8pt, inner sep=0pt] at ([xshift=12pt,yshift=-12pt]frame.north west) {};
    \node[circle, fill=macyellow,minimum size=8pt, inner sep=0pt] at ([xshift=26pt,yshift=-12pt]frame.north west) {};
    \node[circle, fill=macgreen, minimum size=8pt, inner sep=0pt] at ([xshift=40pt,yshift=-12pt]frame.north west) {};
  },
  listing only,
  listing options={
    style=atom-one-dark,
    aboveskip=0pt,
    belowskip=0pt,
    resetmargins=true, % <-- MAGIC 2: Tells listings to ignore \begin{description} indentation
    xleftmargin=0pt,   % Forces text tight against the left edge of the box
    language=sh,
    #1
  }
}